% ══════════════════════════════════════════════════════════════════════
% Chapter 4: Resonance Activation and Cultural Priming
% ══════════════════════════════════════════════════════════════════════

\chapter{Resonance Activation and Cultural Priming}
\label{ch:activation}

\epigraph{``The question of relevance is the question of what is
activated.''}{Dan Sperber, \textit{Explaining Culture}}

% ── Introduction ──────────────────────────────────────────────────────

The concept of activation---which ideas a signal ``turns on'' in a
culturally constituted mind---lies at the intersection of several
major intellectual traditions.  Sperber and Wilson's \emph{relevance
theory}~\cite{sperber1995} articulates two foundational principles.
The \emph{cognitive principle of relevance} holds that human cognition
is geared toward the maximisation of relevance: cognitive systems
automatically allocate attention to inputs that promise the greatest
cognitive effect for the least processing effort.  The
\emph{communicative principle of relevance} adds that every ostensive
stimulus carries a presumption of its own optimal relevance---that the
speaker has made it worth the hearer's effort to
process.\footnote{Sperber and Wilson distinguish ``relevance'' from
mere ``informativeness.''  A stimulus is relevant not because it
carries a large quantity of information, but because it connects
efficiently with the hearer's existing cognitive environment---what we
formalise here as the repertoire.}  IDT's resonance relation gives
these principles a precise algebraic form: $r \res s$ holds exactly
when the signal~$s$ is relevant to the schema~$r$.

Cognitive linguistics provides a complementary perspective.  Lakoff's
\emph{Women, Fire, and Dangerous Things} (1987)~\cite{lakoff1987}
demonstrated that human categorisation is not a matter of classical
necessary-and-sufficient conditions but of radial categories organised
around prototypes and linked by family resemblance.  The schemas that
a signal activates are not an arbitrary set; they form a structured
neighbourhood in conceptual space, held together by what Lakoff calls
\emph{idealized cognitive models} (ICMs).  Fillmore's \emph{frame
semantics}~\cite{fillmore1982} makes a parallel claim at the level of
linguistic meaning: understanding a word requires activating an entire
frame---a structured background of knowledge against which the word
makes sense.\footnote{Fillmore's classic example: the word
``\emph{bachelor}'' activates a frame that includes the institution of
marriage, typical marital age, and social expectations.  Remove the
frame, and the word loses its meaning.  In IDT terms, the frame is
part of the repertoire, and the word resonates with it.}  The
activated subset of a repertoire (Definition~\ref{def:activated-subset})
is, in Fillmorean terms, the set of frames that a signal evokes.

From experimental psychology, Kahneman's \emph{Thinking, Fast and
Slow} (2011)~\cite{kahneman2011} documents the pervasive role of
\emph{priming}---the phenomenon whereby exposure to one stimulus
influences the response to a subsequent stimulus.  Subjects primed
with words associated with old age walk more slowly; subjects primed
with money-related images become more individualistic.  Kahneman
locates priming in ``System~1,'' the fast, automatic, associative mode
of thought.  IDT's activation framework provides an algebraic account
of exactly these effects: the primer~$s_1$ reshapes the repertoire
(via $\operatorname{interpret}(\Rep, s_1)$), and the subsequent
signal~$s_2$ then activates schemas within this reshaped
landscape.\footnote{The sequential interpretation associativity
theorem (Theorem~\ref{thm:sequential-activation}) shows that this
two-step process is equivalent to a single exposure to the compound
signal $s_1 \compose s_2$.  Priming, in algebraic terms, is
composition.}

What IDT contributes to these diverse literatures is a unified formal
language.  Sperber's relevance, Lakoff's ICMs, Fillmore's frames, and
Kahneman's priming are all, at bottom, claims about \emph{which
elements of a structured repertoire a signal activates}.  By
axiomatising resonance and deriving activation as a consequence, IDT
reveals the common algebraic skeleton beneath these superficially
different theories.

When a signal arrives at a mind, it does not activate the entire
repertoire uniformly.  It selectively \emph{activates} those ideas with
which it resonates---Sperber's (1996) relevance theory~\cite{sperber1996}
formalised.  A Catholic seeing a cross activates christological schemas.
A Hindu seeing the same cross might activate geometric symmetry schemas.
The signal is identical; the activation pattern differs because different
repertoire elements resonate with it.

This chapter builds the formal theory of activation: signals do not just
produce interpretations---they \emph{select which interpretations
matter} through resonance.  We connect this to Bourdieu's
\emph{habitus}~\cite{bourdieu1984}: the culturally acquired dispositions
that determine what is ``relevant'' are precisely the resonance structure
of the repertoire.

\medskip
\noindent\textbf{Chapter overview.}  We define activation
(\S\ref{sec:ch4:activation}), study its interaction with interpretation
(\S\ref{sec:ch4:interaction}), analyse sequential (double) activation
(\S\ref{sec:ch4:double}), establish depth bounds
(\S\ref{sec:ch4:depth-bounds}), and trace resonance chains
(\S\ref{sec:ch4:chains}).

\subsection*{Historical Context: From Associationism to Relevance Theory}

The idea that mental contents activate one another has a long
pedigree.  British associationism, from David Hume's \emph{Treatise of
Human Nature} (1739) through David Hartley's \emph{Observations on
Man} (1749), proposed that ideas are linked by resemblance, contiguity,
and cause-and-effect, and that the activation of one idea
automatically triggers associated ideas.  This philosophical tradition
was given cognitive-scientific form by Collins and
Loftus~\cite{collins1975}, whose \emph{spreading activation} model
posited that concepts are stored in a semantic network and that
activating one node sends activation spreading along weighted links to
neighbouring nodes.  The model explained priming effects, category
verification times, and typicality gradients.

Anthropologists have long employed activation concepts, if not always
under that name.  Malinowski's ``context of
situation''~\cite{malinowski1923} was an early recognition that the
meaning of an utterance depends on which aspects of the surrounding
context are cognitively active for the participants.  Bateson's theory
of ``frames''~\cite{bateson1972} proposed that metacommunicative
signals (``this is play'') activate interpretive frames that transform
the meaning of all subsequent signals.  Goffman~\cite{goffman1974}
extended Bateson's frame analysis into a comprehensive sociology of
everyday interaction, showing how \emph{frame shifts}---changes in
which schemas are active---are the fundamental unit of social
micro-dynamics.

IDT inherits and formalises this lineage.  The resonance
relation~$\res$ replaces the informal notion of ``associative link'';
the activated subset replaces the vague concept of ``what is salient'';
and the composition operation~$\compose$ replaces the hand-waving
appeal to ``contextual influence.''  The theorems that follow make
precise the claims that associationists, relevance theorists, and
frame analysts have long advanced informally.

% ── §1. Activation ───────────────────────────────────────────────────
\section{Activation}
\label{sec:ch4:activation}

\begin{definition}[Activation]
\label{def:activation}
An idea~$r$ is \emph{activated} by signal~$s$ if they resonate:
$\operatorname{activated}(r, s) \iff r \res s$.
Activation is the formal counterpart of Sperber's ``relevance'': a
signal is relevant to a schema if and only if it resonates with it.
\end{definition}

\begin{definition}[Activated Subset]
\label{def:activated-subset}
The \emph{activated subset} of repertoire~$\Rep$ under selection
predicate $\sigma$ is $\operatorname{activatedSubset}(\Rep, \sigma)
= \operatorname{filter}(\sigma, \Rep)$.
\end{definition}

\begin{definition}[Activated Interpretations]
\label{def:activated-interp}
The \emph{activated interpretations} compose only the selected elements
with the signal:
\[
  \operatorname{activatedInterpretations}(\Rep, s, \sigma)
  = \operatorname{map}(\lambda r.\; r \compose s,\;
    \operatorname{filter}(\sigma, \Rep)).
\]
This models selective attention: only relevant schemas participate in
interpretation.
\end{definition}

\begin{theorem}[Self-Activation]
\label{thm:self-activation}
Every idea activates itself: $\operatorname{activated}(r, r)$ for all~$r$.
\end{theorem}

\begin{proof}[Proof sketch]
By reflexivity of resonance.  See \texttt{self\_activation}.
\end{proof}

\begin{remark}
The formal basis of self-recognition.  Familiar symbols are immediately
meaningful because they resonate with what is already there.
\end{remark}

\begin{theorem}[Activation Symmetry]
\label{thm:activation-symm}
If $\operatorname{activated}(r, s)$, then $\operatorname{activated}(s, r)$.
\end{theorem}

\begin{proof}[Proof sketch]
By symmetry of resonance.  See \texttt{activation\_symm}.
\end{proof}

\begin{remark}
Relevance is reciprocal.  If a Christian schema is activated by a cross,
then the cross-schema would be activated by a Christian symbol.
Cultural exchange is inherently bidirectional.
\end{remark}

\begin{theorem}[Void Self-Activates]
\label{thm:void-self-activates}
$\operatorname{activated}(\void, \void)$.
\end{theorem}

\begin{proof}[Proof sketch]
By reflexivity: $\void \res \void$.
See \texttt{void\_self\_activates}.
\end{proof}

\begin{theorem}[Composition Preserves Activation (Right)]
\label{thm:activation-compose-right}
If $\operatorname{activated}(r, s)$, then
$\operatorname{activated}(r \compose c,\; s \compose c)$.
\end{theorem}

\begin{proof}[Proof sketch]
By resonance-compose-right.  See \texttt{activation\_compose\_right}.
\end{proof}

\begin{theorem}[Composition Preserves Activation (Left)]
\label{thm:activation-compose-left}
If $\operatorname{activated}(r, s)$, then
$\operatorname{activated}(c \compose r,\; c \compose s)$.
\end{theorem}

\begin{proof}[Proof sketch]
By resonance-compose-left.  See \texttt{activation\_compose\_left}.
\end{proof}

\begin{remark}
A ritual frame (opening prayer, spatial arrangement) that precedes two
resonant symbols does not destroy their mutual activation.  Context
preserves relevance.
\end{remark}

\begin{example}[The Cross as Multi-Activation Symbol]
\label{ex:cross-multi-activation}
The symbol of the cross provides a vivid illustration of how the same
signal activates radically different schemas depending on the
repertoire it encounters.  For a devout Catholic, the cross activates
a dense cluster of christological schemas: the crucifixion narrative,
the theology of redemption through suffering, the eucharistic
sacrifice, Marian devotion at the foot of the cross, and the
eschatological promise of resurrection.  These schemas are mutually
resonant, forming what Lakoff would call an idealized cognitive model
of salvific suffering.

For a mathematician, the identical visual form activates an entirely
disjoint set of schemas: Cartesian coordinate axes, the addition
operator in elementary arithmetic, the cross product of vectors, or
the symbol for the direct product of groups.  The resonance relation
$r \res s$ holds for a different subset of the repertoire, and the
activated interpretations accordingly diverge.

For a Red Cross worker, the cross activates schemas of humanitarian
aid, political neutrality, the Geneva Conventions, and emergency field
medicine.  The symbol's meaning is neither theological nor
mathematical but \emph{institutional}---it indexes a framework of
international law and organisational practice.

In each case, the signal is identical; what differs is the
repertoire~$\Rep$ and hence the activated subset.  The composition
preservation theorems (Theorems~\ref{thm:activation-compose-right}
and~\ref{thm:activation-compose-left}) guarantee that adding
contextual framing to the cross preserves these divergent activation
patterns rather than collapsing them.
\end{example}

\begin{example}[Priming in Advertising]
\label{ex:advertising-priming}
Consider a television advertisement for a luxury automobile.  Before
the car appears on screen, the viewer is shown a sequence of images:
a private jet landing on a sun-drenched runway, a penthouse overlooking
a glittering skyline, a well-dressed figure signing documents in a
boardroom of polished mahogany.  These images are not accidental; they
are primers.  Each image activates schemas of wealth, success,
aspiration, and social status.  By the time the automobile appears,
the viewer's repertoire has been reshaped: the schemas most recently
activated are those associated with luxury and achievement.  The car
is then interpreted not merely as a mode of transportation but as a
signifier of the lifestyle depicted in the preceding images.

The sequential interpretation associativity theorem
(Theorem~\ref{thm:sequential-activation}) explains why this works:
the priming sequence followed by the product reveal is algebraically
equivalent to a single compound signal that fuses aspiration-imagery
with automobile-imagery.  The advertiser's art lies in choosing primers
whose composition with the product signal maximises the desired
activation pattern.
\end{example}

\subsection*{Discussion: Activation as Cultural Attention}

The activation framework developed so far formalises what
anthropologists have long called \emph{cultural salience}---the
culturally variable tendency to notice, attend to, and find meaningful
certain features of the environment while ignoring others.  Mary
Douglas, in \emph{Purity and Danger} (1966)~\cite{douglas1966},
argued that every culture imposes a classification system that directs
attention to particular boundaries (pure/impure, sacred/profane,
inside/outside) and renders invisible everything that does not fit the
classificatory grid.  In IDT terms, Douglas's cultural classifications
are encoded in the repertoire~$\Rep$, and the boundaries she identifies
are precisely the regions of high resonance density: signals that fall
near a classificatory boundary activate many schemas simultaneously,
producing the sense of heightened significance---or danger---that
Douglas documented.

Conversely, signals that fall outside all cultural categories activate
nothing (the null activation theorem, Theorem~\ref{thm:null-activation})
and are experienced as meaningless noise.  Douglas's famous ``matter
out of place''---dirt as the by-product of classification---is, in
formal terms, a signal whose activated subset is empty relative to the
prevailing repertoire.  What counts as dirt, disorder, or anomaly is
not an objective property of the signal but a consequence of the
repertoire's resonance structure.

% ── §2. Activation and Interpretation ────────────────────────────────
\section{Activation and Interpretation Interaction}
\label{sec:ch4:interaction}

\begin{theorem}[Activated Interpretations Are a Sublist]
\label{thm:activated-sublist}
\[
  |\operatorname{activatedInterpretations}(\Rep, s, \sigma)|
  \leq |\operatorname{interpret}(\Rep, s)|.
\]
\end{theorem}

\begin{proof}[Proof sketch]
Filtering can only reduce list length; mapping preserves it.
See \texttt{activated\_interp\_sublist}.
\end{proof}

\begin{remark}
Attention is a filter, not a generator.  Selective attention can only
\emph{reduce} the set of interpretations, never add new ones.
\end{remark}

\begin{theorem}[Universal Activation = Full Interpretation]
\label{thm:universal-activation}
If $\sigma(r) = \mathtt{true}$ for all~$r$, then
$\operatorname{activatedInterpretations}(\Rep, s, \sigma)
 = \operatorname{interpret}(\Rep, s)$.
\end{theorem}

\begin{proof}[Proof sketch]
Filtering with the constant-true predicate is the identity.
See \texttt{universal\_activation}.
\end{proof}

\begin{remark}
A mind with zero attentional filters (the ``open mind'' ideal) interprets
everything.  Pure receptivity means every schema participates.
\end{remark}

\begin{theorem}[Null Activation = Empty Interpretation]
\label{thm:null-activation}
If $\sigma(r) = \mathtt{false}$ for all~$r$, then
$\operatorname{activatedInterpretations}(\Rep, s, \sigma) = []$.
\end{theorem}

\begin{proof}[Proof sketch]
Filtering with the constant-false predicate yields the empty list.
See \texttt{null\_activation}.
\end{proof}

\begin{remark}
Total cognitive shutdown---when nothing in your repertoire is deemed
relevant---produces zero interpretations.  This models culture shock or
the ``deer in headlights'' response to incomprehensible stimuli.
\end{remark}

\begin{example}[Culture Shock as Null Activation]
\label{ex:culture-shock}
An American businessman arrives in a remote village in the Papua New
Guinea highlands.  He is led to a clearing where men in elaborate
feathered headdresses are engaged in a ceremonial exchange of pearl
shells, pigs, and sweet potatoes, accompanied by stylised oratory,
rhythmic stamping, and the blowing of bamboo flutes.  Nothing in his
cognitive repertoire---business negotiation schemas, American social
etiquette, Western aesthetic categories---resonates with what he
observes.  The activation predicate $\sigma(r)$ returns
$\mathtt{false}$ for every element $r \in \Rep$.  The result, by
Theorem~\ref{thm:null-activation}, is an empty list of
interpretations: complete cognitive blankness.

This is the formal model of what Oberg (1960)~\cite{oberg1960} called
\emph{culture shock}---not merely unfamiliarity but the total failure
of one's interpretive apparatus.  The remedy, in IDT terms, is
repertoire expansion: by acquiring new schemas (through language
learning, participant observation, or explicit instruction), the
traveller gradually populates $\Rep$ with elements that resonate with
local signals, restoring the capacity for interpretation.
\end{example}

\begin{remark}[Connection to Artificial Intelligence: Attention Mechanisms]
\label{rem:attention-ai}
The activation framework bears a striking structural parallel with the
\emph{attention mechanism} in transformer neural
networks~\cite{vaswani2017}.  In a transformer, each input token is
projected into query, key, and value vectors; the attention score
between a query and a key determines how much the corresponding value
contributes to the output.  In IDT, a signal~$s$ plays the role of
the query, the repertoire elements~$r \in \Rep$ play the role of keys,
and the interpretations $r \compose s$ play the role of values.  The
resonance predicate $r \res s$ is a binary attention score:
it selects which values contribute and which are suppressed.

This analogy is more than metaphorical.  The activated interpretations
definition (Definition~\ref{def:activated-interp}) is a discrete,
binary version of the attention-weighted sum that transformers
compute.  The null activation theorem
(Theorem~\ref{thm:null-activation}) corresponds to the case where all
attention scores are zero---the model produces no output.  The
universal activation theorem
(Theorem~\ref{thm:universal-activation}) corresponds to uniform
attention---every key-value pair contributes equally.  Future work
might extend IDT's binary resonance to graded resonance, bringing the
formal framework even closer to the continuous attention scores used
in neural architectures.
\end{remark}

% ── §3. Double Activation ────────────────────────────────────────────
\section{Double Activation: Sequential Signals}
\label{sec:ch4:double}

\begin{theorem}[Sequential Interpretation Associativity]
\label{thm:sequential-activation}
\[
  \operatorname{interpret}(\operatorname{interpret}(\Rep, s_1), s_2)
  = \operatorname{interpret}(\Rep, s_1 \compose s_2).
\]
\end{theorem}

\begin{proof}[Proof sketch]
By functoriality of map and associativity of composition.
See \texttt{sequential\_activation}.
\end{proof}

\begin{remark}
Cultural priming works because sequential exposure to signals is
equivalent to a single compound signal.  Ritual sequences (opening chant
$\to$ sacrifice $\to$ closing prayer) produce the same interpretations
as a single complex ritual.
\end{remark}

\begin{theorem}[Resonant Sequential Identity]
\label{thm:resonant-sequential}
For any background~$r$ and signals $s_1, s_2$:
\[
  r \compose (s_1 \compose s_2) = (r \compose s_1) \compose s_2.
\]
\end{theorem}

\begin{proof}[Proof sketch]
By associativity.  See \texttt{resonant\_sequential\_identity}.
\end{proof}

\begin{theorem}[Void Priming is Identity]
\label{thm:void-priming}
\[
  \operatorname{interpret}(\operatorname{interpret}(\Rep, \void), s)
  = \operatorname{interpret}(\Rep, s).
\]
\end{theorem}

\begin{proof}[Proof sketch]
$\operatorname{interpret}(\Rep, \void) = \Rep$ by silence transparency,
so the second interpretation acts on the original repertoire.
See \texttt{void\_priming}.
\end{proof}

\begin{remark}
A ritual preceded by silence produces the same interpretations as the
ritual alone.  The ``pregnant pause'' before a sermon adds drama but not
information.
\end{remark}

\begin{theorem}[Post-Void is Identity]
\label{thm:post-void}
\[
  \operatorname{interpret}(\operatorname{interpret}(\Rep, s), \void)
  = \operatorname{interpret}(\Rep, s).
\]
\end{theorem}

\begin{proof}[Proof sketch]
By silence transparency applied to $\operatorname{interpret}(\Rep, s)$.
See \texttt{post\_void}.
\end{proof}

\subsection*{Discussion: Sequential Activation and Ritual Design}

The theorems of this section reveal that the careful ordering of
ritual elements is not arbitrary but reflects the algebraic structure
of composition.  Ritual designers---priests, shamans, liturgists,
and, in the modern world, advertisers and political
speechwriters---exploit sequential activation to shape the
interpretive landscape of their audiences.  The sequential
interpretation associativity theorem
(Theorem~\ref{thm:sequential-activation}) shows that the effect of a
ritual sequence can be analysed as a single compound signal, while
the void priming theorem (Theorem~\ref{thm:void-priming}) and the
post-void theorem (Theorem~\ref{thm:post-void}) establish that only
\emph{substantive} primers---those carrying non-trivial content---alter
the interpretive outcome.

\begin{example}[The Seder as Sequential Activation]
\label{ex:seder}
The Jewish Passover Seder follows a precise fifteen-step sequence:
\emph{Kadesh} (sanctification over wine), \emph{Urchatz} (hand
washing), \emph{Karpas} (dipping greens in salt water),
\emph{Yachatz} (breaking the middle matzah), \emph{Maggid}
(narrating the Exodus), and so on through \emph{Nirtzah} (acceptance).
Each step activates specific schemas---liberation, slavery, divine
intervention, communal identity, agricultural renewal---that prime
the participant for the next step.  The dipping of bitter herbs
(\emph{Maror}) in \emph{charoset} activates schemas of suffering and
sweetness simultaneously, and this compound activation primes the
participant for the meal (\emph{Shulchan Orech}), which is then
interpreted not merely as eating but as a re-enactment of the
celebratory feast of the freed Israelites.

The sequential associativity theorem tells us that this entire
fifteen-step cascade is algebraically equivalent to a single compound
signal $s_1 \compose s_2 \compose \cdots \compose s_{15}$.  But the
pedagogical genius of the Seder lies in \emph{decomposing} this
compound signal into steps, each of which is individually
interpretable and each of which prepares the cognitive ground for its
successor.
\end{example}

\begin{remark}[Connection to Pedagogy: Scaffolding]
\label{rem:scaffolding}
Vygotsky's concept of the \emph{zone of proximal development}
(ZPD)~\cite{vygotsky1978} and the pedagogical technique of
\emph{scaffolding} (Wood, Bruner, and Ross, 1976) describe the
practice of providing learners with structured support that enables
them to achieve tasks just beyond their current independent
capability.  The void priming theorem
(Theorem~\ref{thm:void-priming}) provides a formal criterion for
effective scaffolding: a preliminary lesson that activates nothing
new---that is, a void primer---adds no interpretive value.  Effective
scaffolding requires \emph{non-void primers} that genuinely activate
relevant schemas, reshaping the learner's repertoire so that the
target material falls within the activated subset.

The post-void theorem (Theorem~\ref{thm:post-void}) adds a
complementary lesson: trailing silence---review sessions that merely
repeat what has already been activated without introducing new
content---likewise adds nothing to the interpretive outcome.  The
implication for instructional design is that both the preamble and the
coda of a lesson must carry substantive content; purely ceremonial
framing is algebraically inert.
\end{remark}

% ── §4. Activation Depth Bounds ──────────────────────────────────────
\section{Activation Depth Bounds}
\label{sec:ch4:depth-bounds}

\begin{theorem}[Activated Interpretation Depth Bound]
\label{thm:activated-depth}
\[
  \operatorname{maxInterpDepth}(\operatorname{activatedInterpretations}(\Rep, s, \sigma))
  \leq \operatorname{maxInterpDepth}(\Rep) + \depth(s).
\]
\end{theorem}

\begin{proof}[Proof sketch]
By induction on $|\Rep|$.  Filtering can only select elements whose
depth is bounded by $\operatorname{maxInterpDepth}(\Rep)$; composition
adds at most $\depth(s)$.
See \texttt{activated\_depth\_bound}.
\end{proof}

\begin{remark}
Selective attention does not change the depth bound.  Relevance-filtering
is a horizontal cut, not a vertical one.
\end{remark}

\begin{theorem}[Activation-Filtered Count Bound]
\label{thm:activated-count}
\[
  |\operatorname{activatedInterpretations}(\Rep, s, \sigma)|
  \leq |\Rep|.
\]
\end{theorem}

\begin{proof}[Proof sketch]
By the length bound on filtered lists.
See \texttt{activated\_count\_bound}.
\end{proof}

\begin{remark}
You cannot hallucinate more schemas than you actually have, no matter how
stimulating the signal.
\end{remark}

\begin{example}[Selective Attention in Ethnography]
\label{ex:ethnographic-attention}
When Clifford Geertz observed the Balinese cockfight~\cite{geertz1973},
his ethnographic training activated schemas that ordinary Balinese
spectators---or casual Western tourists---lacked.  Where a tourist
might activate only schemas of animal cruelty or gambling, Geertz's
repertoire contained schemas of status rivalry, symbolic violence,
Weberian stratification, and Benthamite ``deep play.''  His trained
repertoire therefore produced a richer set of activated
interpretations.  But Theorem~\ref{thm:activated-count} guarantees
that even Geertz could not produce more interpretations than the
cardinality of his repertoire: $|\operatorname{activatedInterpretations}
(\Rep_{\text{Geertz}}, s_{\text{cockfight}}, \sigma)| \leq
|\Rep_{\text{Geertz}}|$.  The ethnographer's advantage is not
unbounded creativity but a \emph{larger and more finely differentiated
repertoire} through which more signals resonate.
\end{example}

\begin{remark}[Depth Bounds and Information Overload]
\label{rem:information-overload}
The activated interpretation depth bound
(Theorem~\ref{thm:activated-depth}) provides a formal explanation for
why information overload, though experientially overwhelming, is
\emph{cognitively bounded}.  In an age of constant digital
stimulation---news feeds, social media notifications, algorithmic
content streams---each individual signal can activate at most
$|\Rep|$ schemas, and the depth of any resulting interpretation is
bounded by $\operatorname{maxInterpDepth}(\Rep) + \depth(s)$.  The
subjective experience of overload arises not from unbounded depth of
processing but from the \emph{rapid succession} of signals, each
making finite but non-negligible demands on interpretive resources.

The so-called ``attention economy''---the competition among signals
for limited cognitive bandwidth---is thus formally bounded by
repertoire structure.  No amount of algorithmic optimisation can force
a mind to produce interpretations deeper or more numerous than its
repertoire permits.  The defense against information overload is not
faster processing but \emph{repertoire curation}: the deliberate
cultivation of schemas that resonate with signals one judges
important, and the pruning of schemas that resonate with noise.
\end{remark}

% ── §5. Resonance Chains ─────────────────────────────────────────────
\section{Resonance Chains and Activation Propagation}
\label{sec:ch4:chains}

\begin{theorem}[Resonant Backgrounds Produce Resonant Activations]
\label{thm:resonant-activations}
If $r_1 \res r_2$, then
$r_1 \compose s \;\res\; r_2 \compose s$.
\end{theorem}

\begin{proof}[Proof sketch]
By resonance-compose-right.
See \texttt{resonant\_backgrounds\_resonant\_activations}.
\end{proof}

\begin{remark}
Culturally similar people who are primed by the same stimulus produce
resonant responses.  This is Durkheim's collective effervescence: a
crowd sharing cultural schemas, exposed to the same ritual stimulus,
produces a harmonious collective response~\cite{durkheim1912}.
\end{remark}

\begin{theorem}[Activation Under Composition]
\label{thm:activation-composition}
If $\operatorname{activated}(r, s_1)$ and
$\operatorname{activated}(s_1, s_2)$, then
$\operatorname{activated}(r \compose s_1,\; s_1 \compose s_2)$.
\end{theorem}

\begin{proof}[Proof sketch]
By the full resonance-compose lemma.
See \texttt{activation\_composition\_chain}.
\end{proof}

\begin{remark}
Resonance ``propagates'' through compositional chains.  Multi-layered
rituals create cascading activations through the audience's schemas.
\end{remark}

\begin{theorem}[Void Is Universally Self-Activatable]
\label{thm:void-self-activatable}
$\operatorname{activated}(\void, \void)$.
\end{theorem}

\begin{proof}[Proof sketch]
By reflexivity of resonance.
See \texttt{void\_universally\_self\_activatable}.
\end{proof}

\begin{theorem}[Double Signal Depth Bound]
\label{thm:double-signal-depth}
\[
  \depth(r_2 \compose (r_1 \compose s))
  \leq \depth(r_2) + \depth(r_1) + \depth(s).
\]
\end{theorem}

\begin{proof}[Proof sketch]
Apply depth-compose-le twice.  See \texttt{double\_signal\_depth\_bound}.
\end{proof}

\begin{remark}
Cultural complexity does not compound super-linearly through chains of
interpretation.  A three-layer cultural transmission has bounded total
complexity.
\end{remark}

\begin{example}[Propaganda as Activation Chain]
\label{ex:propaganda}
The propaganda apparatus of totalitarian regimes provides a dark but
instructive illustration of activation chains.  Joseph Goebbels's
communication strategy in Nazi Germany involved the careful sequencing
of messages: initial signals activated schemas of national humiliation
(the Treaty of Versailles, economic hardship), which primed the
population for subsequent signals activating schemas of ethnic
scapegoating, which in turn primed for signals activating schemas of
militaristic renewal.  Each stage in the sequence depended on the
activation patterns established by its predecessors.

The activation composition theorem
(Theorem~\ref{thm:activation-composition}) shows that this cascade of
resonances propagates through compositional chains: if
$\operatorname{activated}(r, s_1)$ and
$\operatorname{activated}(s_1, s_2)$, then the compositions
$r \compose s_1$ and $s_1 \compose s_2$ are themselves mutually
activated.  But the double signal depth bound
(Theorem~\ref{thm:double-signal-depth}) provides a crucial
counterpoint: the total interpretive depth of such a chain is bounded
by the sum of the individual depths, not their product.  Propaganda
is dangerous, but its cognitive reach is formally limited by the
structure of the repertoires it targets.
\end{example}

\begin{remark}[Open Questions]
\label{rem:open-questions}
The activation framework developed in this chapter treats resonance as
a binary relation: a signal either resonates with a schema or it does
not.  Several natural extensions suggest themselves.

First, can activation be made \emph{quantitative} beyond the binary
resonates/does-not-resonate dichotomy?  A graded resonance
function $\rho(r, s) \in [0, 1]$ would allow modelling of
\emph{partial activation} and \emph{weak relevance}---phenomena that
are ubiquitous in everyday cognition.  A faint church bell in the
distance partially activates religious schemas without fully engaging
them; a foreign word that \emph{sounds like} a familiar word in one's
own language triggers tentative, low-confidence activation.  Graded
resonance connects naturally to fuzzy set theory~\cite{zadeh1965}
and to the graded category membership documented in cognitive
linguistics~\cite{rosch1975}.

Second, the current framework treats the repertoire as static during
a single activation event.  But activation itself may \emph{modify}
the repertoire: schemas that are frequently activated may become more
accessible (the ``frequency effect'' in psycholinguistics), while
schemas that are never activated may atrophy (the ``use it or lose it''
principle).  A dynamic activation theory would require enriching the
algebraic structure with a feedback mechanism from activation to
repertoire composition.

Third, the interaction between activation and social structure remains
largely unexplored.  When two minds with different repertoires are
exposed to the same signal, their divergent activation patterns may
generate communicative friction---or, under favourable conditions,
productive misunderstanding that drives cultural innovation.
Formalising this social dimension of activation is a task for future
chapters.
\end{remark}

% ── Summary ──────────────────────────────────────────────────────────
\section*{Chapter Summary}

This chapter formalised the theory of \emph{resonance activation}:
signals selectively activate repertoire elements through resonance,
producing filtered interpretations.  Self-activation
(Theorem~\ref{thm:self-activation}) and activation symmetry
(Theorem~\ref{thm:activation-symm}) establish the basic structure.
Universal activation recovers full interpretation
(Theorem~\ref{thm:universal-activation}), while null activation produces
cognitive shutdown (Theorem~\ref{thm:null-activation}).  Sequential
signals compose associatively
(Theorem~\ref{thm:sequential-activation}), and void priming is the
identity (Theorem~\ref{thm:void-priming}).  The activated depth bound
(Theorem~\ref{thm:activated-depth}) shows that selective attention does
not change the depth ceiling.  Resonance propagates through activation
chains (Theorem~\ref{thm:activation-composition}), providing the formal
basis for understanding how multi-layered rituals create cascading
effects through culturally constituted minds.

All eighteen theorems are verified in
\texttt{FormalAnthropology/IDT\_Signal\_Ch04.lean}.
